\documentclass[11pt]{article}
\usepackage[a4paper,margin=25mm]{geometry}
\usepackage[T1]{fontenc}
\usepackage{lmodern}
\usepackage{booktabs}
\usepackage{hyperref}
\usepackage{xcolor}
\usepackage{listings}
\lstset{basicstyle=\ttfamily\small, columns=fullflexible, keepspaces=true,
  breaklines=true, frame=single, framerule=0.3pt, rulecolor=\color{gray}}
\newcommand{\code}[1]{\texttt{#1}}
\newcommand{\us}{\textmu s}

\title{\code{runinfo.py}: run conditions, acquisition rate and true trigger rate}
\author{BGOTSCAN}
\date{2026-10-14}

\begin{document}
\maketitle

\section{Purpose}

One table with the conditions of every run of the scan (start time, the
three thermometers, bias, threshold) and what the event bookkeeping of the
file says about the acquisition: how long the 5000 events took, at what rate
they were recorded, how many triggers actually fired on the board meanwhile,
and how the recorded events are spaced in time. It reads only the attributes
and the two small per-event datasets \code{timestamp\_ns} and
\code{trigger\_id} of each file, never the waveforms, so it runs in seconds
over the twelve runs. It is a Python script in \code{src/py/}.

\section{Usage}

\begin{lstlisting}
python3 src/py/runinfo.py [options]
\end{lstlisting}

\begin{center}
\begin{tabular}{l l p{78mm}}
\toprule
option & default & meaning \\
\midrule
\code{--data DIR} & \code{2026\_10\_02\_bgo\_cryo\_tempscan/data\_cryo} (relative to the repository) & directory of the run files; \code{temperatures.txt} is looked for in its parent \\
\code{--runs A B ...} & all \code{*.h5} in the directory & restrict to these runs (names with or without \code{.h5}) \\
\code{--gap MS} & 10 & an interval between consecutive recorded events longer than this (ms) is counted as a readout gap and separates two bursts \\
\code{--save NAME} & off & save the figure as \code{src/py/plots/NAME\_rates.png}; without it the figure is shown on screen \\
\code{--no-figure} & off & table only \\
\bottomrule
\end{tabular}
\end{center}

\section{Inputs}

\paragraph{Run name.} \code{TTTK\_BBVV\_thNN}: nominal crystal temperature in
K, SiPM bias in V with \code{d} for the decimal point (\code{39d5V} is
39.5~V), trigger threshold in ADC counts.

\paragraph{\code{temperatures.txt}.} Laura's file with three Python lists
(\code{names}, \code{T1}, \code{T2}, \code{T3}) in the same order. T1 is on
the aluminium case, T2 under it, T3 on the back face of the crystal. The
script matches runs to this file by name; a run absent from it gets blanks.

\paragraph{File attributes used.} \code{start\_time}, \code{events},
\code{itl\_majority}, \code{threshold\_mode}, \code{pretrigger},
\code{record\_length}, \code{tsample\_ns}.

\paragraph{\code{timestamp\_ns}.} The board time of each recorded event in
ns, from the digitizer clock.

\paragraph{\code{trigger\_id}.} The board's trigger counter for each
recorded event. The board increments it for \emph{every} trigger of the
ITL logic, whether or not the event was read out, and it is a 24-bit
counter that wraps at $2^{24} = 16\,777\,216$. The difference of consecutive
recorded values is therefore the number of triggers that fired between the
two events, including the recorded one, modulo $2^{24}$.

\section{Definitions}

\paragraph{Duration.} $D = (t_{N-1} - t_0)$ in seconds, from the first to
the last recorded event.

\paragraph{Recorded rate.} $r_{\rm rec} = (N-1)/D$ with $N$ the number of
events in the file.

\paragraph{True trigger rate.} The number of board triggers over the run,
\[
  n_{\rm trig} = (\code{trigger\_id}_{N-1} - \code{trigger\_id}_0) \bmod 2^{24},
\]
divided by $D$: $r_{\rm trig} = n_{\rm trig}/D$. The modulo handles a single
wrap of the counter during the run; the counter can only wrap once in
40~s (it would need more than 400~kHz). The fraction of triggers that were
recorded is $f = (N-1)/n_{\rm trig}$.

\paragraph{Bursts and gaps.} The intervals $\Delta t_i = t_{i+1} - t_i$
between consecutive recorded events are split at \code{--gap}: intervals
longer than it are readout gaps, the runs of events between gaps are
bursts. The table gives the number of bursts, the median burst size (events)
and the median gap and in-burst interval (ms). A readout that keeps up with
the trigger has no gaps; a saturated one shows a fixed burst size (the
board's buffer) and a fixed gap (the transfer time).

\section{Output}

\paragraph{Table.} One line per run, sorted by \code{start\_time}, with:
run name, start time (clock time), T1, T2, T3 (K), bias (V), threshold (ADC),
majority, duration (s), recorded rate (Hz), true trigger rate (Hz), fraction
recorded (\%), number of bursts, median events per burst, median gap (ms),
median in-burst interval (ms). The last line recalls the common settings
(events, record length, pre-trigger, sampling).

\paragraph{Figure (\code{NAME\_rates.png}).} Two panels, one $y$ axis each,
logarithmic in rate.
\begin{itemize}
\item Left: true trigger rate versus bias. Runs are grouped by nominal
  temperature; the groups taken at one setpoint (184--187~K and 195--197~K)
  are joined by a line, the single runs (128, 153, 164~K) are isolated
  points. Each point is labelled with its T3. The recorded rate is drawn
  as a grey horizontal line: it is the same in all runs and is the readout
  limit, not a physical rate.
\item Right: true trigger rate versus T3, each point labelled with its
  bias. The runs taken at ``constant overvoltage'' (bias raised with
  temperature) are joined by a line.
\end{itemize}

\section{What the first run of the script showed (2026-10-14)}

All twelve runs record 5000 events in 40.7~s, i.e.\ 123~Hz regardless of
threshold, bias and temperature; the events come in bursts of about three,
1.3~ms apart, separated by gaps of 22.8~ms. That is a readout saturated at
about 47~MB/s. The trigger counter gives true rates from 1.9~kHz (128~K,
37~V) to 225~kHz (195~K, 41~V), rising steeply with bias at fixed
temperature and with temperature at roughly constant overvoltage: the
trigger is dominated by dark-count coincidences, and the recorded events
are a sample of 0.05\% to 7\% of the triggers. The scintillation events
are a minority of the sample at the warm end, and an event selection that
finds the pulse is needed before any light-yield measurement.

\section{Examples}

\begin{lstlisting}
# table and figure on screen
python3 src/py/runinfo.py

# table only
python3 src/py/runinfo.py --no-figure

# figure saved to src/py/plots/scan_rates.png
python3 src/py/runinfo.py --save scan
\end{lstlisting}

\end{document}
